% GENERATED FILE. Edit canonical lessons, then run npm run generate:books.
% canonical-source-hash: fnv1a64:d09247dc5fbac6d3
% canonical-lessons: RU-C374-rasschityvat, RU-C374-ozhidat, RU-C374-ozhidanie, RU-C374-udavatsya, RU-C374-poluchatsya

\chapter{Expectations}
\label{ch:ru-expectations}
\hlchaptermodality{374}

\begin{chapteropening}
\textbf{By the end of this chapter:} \emph{I can say what I count on and what I expect, and whether I manage to do something or it works out.}

Complete the last lesson of chapter 374: I can say what I count on and what I expect, and whether I manage to do something or it works out.
\end{chapteropening}

\section[\texorpdfstring{\ru{рассчитывать} (rasschítyvat)}{рассчитывать (rasschítyvat)}]{\ru{рассчитывать} (rasschítyvat) — to count on, to rely on; to calculate}

\label{lesson:RU-C374-rasschityvat}



\begin{quote}
Before the new one: say the Russian for to threaten, then the Russian for to prevent.
\end{quote}

\subsection*{You'll want to know: \ru{рассчитывать}}
\textbf{\ru{рассчитывать}} (\emph{rasschítyvat}) — \textquotedblleft{}to count on, to rely on; to calculate\textquotedblright{}.

\begin{quote}
\ru{Мы} \ru{рассчитываем} \ru{на} \ru{твою} \ru{помощь}.
\end{quote}

\begin{quote}
\emph{My rasschítyvayem na tvoyú pómoshch.}
\end{quote}

\begin{quote}
We are counting on your help.
\end{quote}

Imperfective; perfective \emph{\ru{рассчитать}} (to calculate). \emph{\ru{Рассчитывать} \ru{на}} + accusative — to count on.

\begin{cousinweb}
\textbf{\ru{рассчитывать}} is built from \emph{\ru{рас}-} and \emph{\ru{считать}} \textquotedblleft{}to count\textquotedblright{}, on the root of \emph{\ru{число}} \textquotedblleft{}number\textquotedblright{}, Proto-Slavic \emph{*čislo} (EtymDB). To count on someone was first to reckon them in.
\end{cousinweb}

\begin{grammarlens}[title={What you've built}]
One more word for telling, explaining and describing at B1.
\end{grammarlens}

\subsection*{Your turn}
\begin{itemize}
\raggedright
  \item \emph{Say it:} \emph{rasschítyvat}
  \item \emph{Say it:} \emph{rasschítyvat}, once more
  \item \emph{Say it:} a sentence of your own with \emph{rasschítyvat}
  \item \emph{Recall:} say the Russian for to threaten, then the Russian for to prevent, then say \emph{rasschítyvat} again
\end{itemize}

\begin{tcolorbox}[breakable,colback=teal!4,colframe=teal!35!black,title={Before you move on}]
What does \ru{рассчитывать} mean? (\textquotedblleft{}To count on, to rely on; to calculate\textquotedblright{}.) Say it once more.
\end{tcolorbox}
\section[\texorpdfstring{\ru{ожидать} (ozhidát)}{ожидать (ozhidát)}]{\ru{ожидать} (ozhidát) — to expect, to await}

\label{lesson:RU-C374-ozhidat}



\begin{quote}
Before the new one: say the Russian for to prevent, then the Russian for to count on.
\end{quote}

\subsection*{You'll want to know: \ru{ожидать}}
\textbf{\ru{ожидать}} (\emph{ozhidát}) — \textquotedblleft{}to expect, to await\textquotedblright{}.

\begin{quote}
\ru{Мы} \ru{не} \ru{ожидали} \ru{такого} \ru{результата}.
\end{quote}

\begin{quote}
\emph{My ne ozhidáli takógo rezultáta.}
\end{quote}

\begin{quote}
We did not expect such a result.
\end{quote}

Imperfective. It takes the genitive or the accusative: \emph{\ru{ожидать} \ru{ответа}}. More formal than \emph{\ru{ждать}}.

\begin{cousinweb}
\textbf{\ru{ожидать}} is built from \emph{\ru{о}-} and \emph{-\ru{жидать}}, the long shape of the root of \emph{\ru{ждать}} \textquotedblleft{}to wait\textquotedblright{} (as \emph{-\ru{бирать}} stands beside \emph{\ru{брать}}); \emph{\ru{ждать}} is inherited from Proto-Slavic \emph{*žĭdati} \textquotedblleft{}to wait, to await\textquotedblright{} (EtymDB).
\end{cousinweb}

\begin{grammarlens}[title={What you've built}]
One more word for telling, explaining and describing at B1.
\end{grammarlens}

\subsection*{Your turn}
\begin{itemize}
\raggedright
  \item \emph{Say it:} \emph{ozhidát}
  \item \emph{Say it:} \emph{ozhidát}, once more
  \item \emph{Say it:} a sentence of your own with \emph{ozhidát}
  \item \emph{Recall:} say the Russian for to prevent, then the Russian for to count on, then say \emph{ozhidát} again
\end{itemize}

\begin{tcolorbox}[breakable,colback=teal!4,colframe=teal!35!black,title={Before you move on}]
What does \ru{ожидать} mean? (\textquotedblleft{}To expect, to await\textquotedblright{}.) Say it once more.
\end{tcolorbox}
\section[\texorpdfstring{\ru{ожидание} (ozhidániye)}{ожидание (ozhidániye)}]{\ru{ожидание} (ozhidániye) — an expectation; waiting}

\label{lesson:RU-C374-ozhidanie}



\begin{quote}
Before the new one: say the Russian for to count on, then the Russian for to expect.
\end{quote}

\subsection*{You'll want to know: \ru{ожидание}}
\textbf{\ru{ожидание}} (\emph{ozhidániye}) — \textquotedblleft{}an expectation; waiting\textquotedblright{}.

\begin{quote}
\ru{Результат} \ru{превзошёл} \ru{все} \ru{ожидания}.
\end{quote}

\begin{quote}
\emph{Rezultát prevzoshyól vse ozhidániya.}
\end{quote}

\begin{quote}
The result exceeded all expectations.
\end{quote}

Neuter. \emph{\ru{Зал} \ru{ожидания}} — a waiting room.

\begin{cousinweb}
\textbf{\ru{ожидание}} is built on \emph{\ru{ожидать}} \textquotedblleft{}to expect\textquotedblright{} with the noun suffix \emph{-\ru{ание}}; the root is that of \emph{\ru{ждать}} \textquotedblleft{}to wait\textquotedblright{}.
\end{cousinweb}

\begin{grammarlens}[title={What you've built}]
One more word for telling, explaining and describing at B1.
\end{grammarlens}

\subsection*{Your turn}
\begin{itemize}
\raggedright
  \item \emph{Say it:} \emph{ozhidániye}
  \item \emph{Say it:} \emph{ozhidániye}, once more
  \item \emph{Say it:} a sentence of your own with \emph{ozhidániye}
  \item \emph{Recall:} say the Russian for to count on, then the Russian for to expect, then say \emph{ozhidániye} again
\end{itemize}

\begin{tcolorbox}[breakable,colback=teal!4,colframe=teal!35!black,title={Before you move on}]
What does \ru{ожидание} mean? (\textquotedblleft{}An expectation; waiting\textquotedblright{}.) Say it once more.
\end{tcolorbox}
\section[\texorpdfstring{\ru{удаваться} (udavátsya)}{удаваться (udavátsya)}]{\ru{удаваться} (udavátsya) — to succeed, to manage (impersonal)}

\label{lesson:RU-C374-udavatsya}



\begin{quote}
Before the new one: say the Russian for to expect, then the Russian for an expectation.
\end{quote}

\subsection*{You'll want to know: \ru{удаваться}}
\textbf{\ru{удаваться}} (\emph{udavátsya}) — \textquotedblleft{}to succeed, to manage (impersonal)\textquotedblright{}.

\begin{quote}
\ru{Мне} \ru{не} \ru{удаётся} \ru{выспаться}.
\end{quote}

\begin{quote}
\emph{Mne ne udayótsya výspatsya.}
\end{quote}

\begin{quote}
I cannot manage to get enough sleep.
\end{quote}

Imperfective; present \emph{\ru{удаётся}}; perfective \emph{\ru{удаться}} (\emph{\ru{удастся}}). Impersonal, with the dative: \emph{\ru{мне} \ru{удалось}} — I managed.

\begin{cousinweb}
\textbf{\ru{удаваться}} is built from \emph{\ru{у}-} and \emph{\ru{давать}} \textquotedblleft{}to give\textquotedblright{}, the ongoing partner of \emph{\ru{дать}}: when a thing succeeds, it gives itself to you.
\end{cousinweb}

\begin{grammarlens}[title={What you've built}]
One more word for telling, explaining and describing at B1.
\end{grammarlens}

\subsection*{Your turn}
\begin{itemize}
\raggedright
  \item \emph{Say it:} \emph{udavátsya}
  \item \emph{Say it:} \emph{udavátsya}, once more
  \item \emph{Say it:} a sentence of your own with \emph{udavátsya}
  \item \emph{Recall:} say the Russian for to expect, then the Russian for an expectation, then say \emph{udavátsya} again
\end{itemize}

\begin{tcolorbox}[breakable,colback=teal!4,colframe=teal!35!black,title={Before you move on}]
What does \ru{удаваться} mean? (\textquotedblleft{}To succeed, to manage (impersonal)\textquotedblright{}.) Say it once more.
\end{tcolorbox}
\section[\texorpdfstring{\ru{получаться} (poluchátsya)}{получаться (poluchátsya)}]{\ru{получаться} (poluchátsya) — to work out, to turn out}

\label{lesson:RU-C374-poluchatsya}



\begin{quote}
Before the new one: say the Russian for an expectation, then the Russian for to succeed.
\end{quote}

\subsection*{You'll want to know: \ru{получаться}}
\textbf{\ru{получаться}} (\emph{poluchátsya}) — \textquotedblleft{}to work out, to turn out\textquotedblright{}.

\begin{quote}
\ru{У} \ru{меня} \ru{не} \ru{получается} \ru{это} \ru{упражнение}.
\end{quote}

\begin{quote}
\emph{U menyá ne polucháyetsya éto uprazhnéniye.}
\end{quote}

\begin{quote}
I cannot get this exercise right.
\end{quote}

Imperfective; perfective \emph{\ru{получиться}}. \emph{\ru{У} \ru{меня} \ru{получается}} — I am getting it right. \emph{\ru{Получилось}!} — it worked!

\begin{cousinweb}
\textbf{\ru{получаться}} is the reflexive of \emph{\ru{получать}}, from \emph{\ru{по}-} and \emph{\ru{лучать}}, from \emph{\ru{лучить}} \textquotedblleft{}to come upon\textquotedblright{}, Proto-Slavic \emph{*lučiti}, tentatively \textquotedblleft{}to find, to reach\textquotedblright{} (EtymDB): what you come upon in the end.
\end{cousinweb}

\begin{grammarlens}[title={What you've built}]
One more word for telling, explaining and describing at B1.
\end{grammarlens}

\subsection*{Your turn}
\begin{itemize}
\raggedright
  \item \emph{Say it:} \emph{poluchátsya}
  \item \emph{Say it:} \emph{poluchátsya}, once more
  \item \emph{Say it:} a sentence of your own with \emph{poluchátsya}
  \item \emph{Recall:} say the Russian for an expectation, then the Russian for to succeed, then say \emph{poluchátsya} again
\end{itemize}

\begin{tcolorbox}[breakable,colback=teal!4,colframe=teal!35!black,title={Before you move on}]
What does \ru{получаться} mean? (\textquotedblleft{}To work out, to turn out\textquotedblright{}.) Say it once more.
\end{tcolorbox}
